\documentclass[11pt, a4paper]{article}

% --- Packages ---
\usepackage[a4paper, top=2.5cm, bottom=2.5cm, left=2.5cm, right=2.5cm]{geometry}
\usepackage{amsmath, amssymb}
\usepackage{graphicx}
\usepackage{booktabs}
\usepackage{parskip}
\usepackage{microtype}
\usepackage{fancyhdr}
\usepackage{titlesec}
\usepackage{xcolor}
\usepackage{mdframed}
\usepackage{enumitem}
\usepackage{hyperref}
\usepackage{caption}
\usepackage{float}

\hypersetup{colorlinks=true, linkcolor=black, urlcolor=blue, citecolor=black}

% --- Page style ---
\pagestyle{fancy}
\fancyhf{}
\renewcommand{\headrulewidth}{0.4pt}
\renewcommand{\footrulewidth}{0.4pt}
\fancyhead[L]{\small\textit{Study Notes: C-MAPSS --- Damage Propagation Modeling}}
\fancyhead[R]{\small\thepage}
\fancyfoot[C]{\small\textit{``I'm not stupid, I'm just slow''}}

% --- Section style ---
\titleformat{\section}{\large\bfseries}{}{0em}{}[\titlerule]
\titleformat{\subsection}{\normalsize\bfseries}{}{0em}{}
\titleformat{\subsubsection}{\normalsize\itshape}{}{0em}{}

% --- Highlight boxes ---
\newmdenv[
  backgroundcolor=gray!8,
  linecolor=gray!40,
  linewidth=0.8pt,
  innerleftmargin=10pt,
  innerrightmargin=10pt,
  innertopmargin=8pt,
  innerbottommargin=8pt,
  skipabove=8pt,
  skipbelow=8pt
]{keybox}

\newmdenv[
  backgroundcolor=orange!6,
  linecolor=orange!40,
  linewidth=0.8pt,
  innerleftmargin=10pt,
  innerrightmargin=10pt,
  innertopmargin=8pt,
  innerbottommargin=8pt,
  skipabove=8pt,
  skipbelow=8pt
]{insightbox}

% --- Caption style ---
\captionsetup{font=small, labelfont=bf, skip=4pt}

% -------------------------------------------------------------------
\begin{document}
% -------------------------------------------------------------------

{\Large\bfseries Reading Notes: Damage Propagation Modeling\\
for Aircraft Engine Run-to-Failure Simulation}

\smallskip
{\small Saxena, Goebel, Simon, Eklund \quad --- \quad PHM'08 \quad --- \quad
\textit{In-depth study edition}}

\medskip\hrule

\tableofcontents

\bigskip\hrule\bigskip

% -------------------------------------------------------------------
\section{The Problem Being Solved}
% -------------------------------------------------------------------

The paper opens with a simple and uncomfortable fact: data-driven prognostics
--- the field of predicting when a machine will fail --- cannot progress without
\emph{run-to-failure data}. And run-to-failure data for aircraft engines is
nearly impossible to obtain.

Real engines do not fail on a convenient schedule. Operators who accumulate
failure records keep them private. Fielded systems are often under-instrumented.
The result is that researchers have no shared ground truth against which to
test and compare their algorithms.

\begin{keybox}
\textbf{The core question the paper answers:}
Can we use a physics-based simulator to manufacture believable run-to-failure
data --- good enough to train and benchmark prognostics algorithms?
\end{keybox}

The answer the authors demonstrate is yes. But the answer requires solving
three sub-problems in sequence:

\begin{enumerate}[noitemsep]
  \item Finding a simulator flexible enough to accept injected degradation
  \item Deciding how to model the evolution of that degradation over time
  \item Defining when the simulated engine has ``failed''
\end{enumerate}

Each of the paper's sections addresses one of these in turn.

% -------------------------------------------------------------------
\section{PHM: Where Prognostics Sits}
% -------------------------------------------------------------------

Before going further, it helps to understand what \textbf{PHM} ---
Prognostics and Health Management --- actually is, and where prognostics
fits inside it.

PHM is a five-layer ecosystem, shown in Figure~\ref{fig:phm_stack}.
Each layer answers a more demanding question than the one below it.

\begin{figure}[H]
  \centering
  \includegraphics[width=0.82\textwidth]{01_phm_ecosystem_stack.png}
  \caption{The PHM ecosystem. Prognostics (highlighted) sits at layer 4,
  between diagnostics and operational decision-making. This paper scopes
  exclusively to that layer.}
  \label{fig:phm_stack}
\end{figure}

The paper explicitly defines prognostics as \emph{estimation of remaining
useful life (RUL)} --- not fault detection, not fault identification.
This scoping decision is important: everything in the paper is designed
to support RUL estimation, and nothing else.

\begin{insightbox}
\textbf{Why does the scoping matter?}
A diagnostic algorithm asks ``what is wrong?'' A prognostic algorithm
asks ``how long is left?'' These are different inference problems requiring
different data, different models, and different evaluation metrics. The
C-MAPSS dataset was built for the second question only. Using it for the
first requires caution.
\end{insightbox}

The larger PHM chain beyond prognostics --- maintenance decisions and
operational planning --- is implied by the scoring function in Section~VII
of the paper. The asymmetric penalty (predicting late is worse than
predicting early) only makes sense if we remember that the \emph{purpose}
of RUL prediction is to inform decisions with real safety consequences.

% -------------------------------------------------------------------
\section{The Inverse Problem: Why PHM Is Hard}
% -------------------------------------------------------------------

The most important conceptual framing in the paper is one it never
explicitly names: PHM is an \textbf{inverse problem}.

Figure~\ref{fig:inverse} shows the two directions of the problem.

\begin{figure}[H]
  \centering
  \includegraphics[width=0.88\textwidth]{02_forward_vs_inverse.png}
  \caption{The forward problem (simulation) is easy. The inverse problem
  (PHM inference) is hard. The hidden degradation state $e(t), f(t)$
  is never directly observable --- only its projection through engine
  physics onto sensor signals can be seen.}
  \label{fig:inverse}
\end{figure}

\subsection{The forward problem}

In simulation, the direction is:
\[
  (e(t),\, f(t)) \;\xrightarrow{\text{C-MAPSS physics}}\; \text{sensor outputs}
\]

This is straightforward. You set the health parameters, run the model,
and read off the sensor values. The mapping is deterministic and well-defined.

\subsection{The inverse problem}

In real PHM deployment, the direction is reversed:
\[
  \text{sensor outputs} \;\xrightarrow{?}\; (e(t),\, f(t)) \;\xrightarrow{?}\; \text{RUL}
\]

You observe only the sensor signals. You must infer the hidden degradation
state, and from that infer how much life remains. This is harder for
several reasons:

\begin{itemize}[noitemsep]
  \item The mapping from $(e, f)$ to sensors is \emph{nonlinear}
  \item The controller actively compensates for degradation,
        obscuring it in the sensor record
  \item The true $(e, f)$ values are never revealed to the algorithm
  \item Noise at multiple stages further blurs the signal
\end{itemize}

\begin{insightbox}
\textbf{The deep implication:} the C-MAPSS dataset is designed so that
the algorithm \emph{cannot} solve the inverse problem directly. It must
learn the statistical relationship between sensor trajectories and RUL
from the training data. This is precisely what makes it a meaningful
benchmark.
\end{insightbox}

% -------------------------------------------------------------------
\section{The Tool: C-MAPSS}
% -------------------------------------------------------------------

\subsection{What C-MAPSS is}

C-MAPSS (Commercial Modular Aero-Propulsion System Simulation) is a
NASA MATLAB/Simulink model of a large commercial turbofan engine
(90,000\,lb thrust class). It was the right choice for this work because
it meets the key requirement: it accepts \emph{health parameter inputs}
that allow the user to inject degradation, and it produces \emph{sensor
outputs} that respond accordingly.

Figure~\ref{fig:engine} shows the engine's five rotating modules and
the overall flow path.

\begin{figure}[H]
  \centering
  \includegraphics[width=\textwidth]{03_cmapss_engine_schematic.png}
  \caption{C-MAPSS turbofan engine schematic. Air enters at the fan
  and exits through the nozzle, passing through LPC, HPC, the combustor,
  HPT, and LPT in sequence. The HPC module (amber) is the one degraded
  in the competition dataset. Health modifiers enter each module; 21
  sensor outputs are recorded; hidden margins feed the health index.}
  \label{fig:engine}
\end{figure}

\subsection{Inputs: health parameter modifiers}

C-MAPSS accepts 14 inputs. Thirteen of them are health modifiers ---
efficiency, flow, and pressure-ratio modifiers for each of the five
rotating components. The fourteenth is fuel flow.

The health modifiers are the ``knobs'' of the degradation injection.
By reducing the HPC efficiency modifier and the HPC flow modifier, the
authors simulate a compressor that is progressively losing performance.

\begin{keybox}
\textbf{Important constraint:} for the competition dataset, only HPC
degradation was simulated. Fan, LPC, HPT, and LPT were left at healthy
values. This was an intentional simplification to keep the benchmark
focused.
\end{keybox}

\subsection{Outputs: sensor measurements and margins}

C-MAPSS produces 58 outputs. The authors selected 21 for the dataset,
covering temperatures ($T_2, T_{24}, T_{30}, T_{50}$), pressures
($P_2, P_{15}, P_{30}, P_{s30}$), speeds ($N_f, N_c, NR_f, NR_c$),
flow ratios ($\phi, BPR, farB$), bleed enthalpies, and demanded values.

Four additional outputs --- $T_{48}$ (EGT), $SmFan$, $SmLPC$,
$SmHPC$ --- are used internally to compute the health index but are
\emph{not} provided to competition participants. This is deliberate:
participants must infer health from the 21 observable sensors, not from
the margins directly.

\subsection{Closed-loop operation and the control problem}

A critical detail: the authors ran C-MAPSS in \textbf{closed-loop}
mode, meaning the engine's built-in controller was active during all
simulations. This has a significant consequence.

The controller's job is to maintain commanded fan speed and to prevent
the engine from exceeding design limits. As the HPC degrades, the
controller compensates --- it adjusts fuel flow to maintain thrust.
This means the sensor signals do not directly reveal the degradation.
Instead, they show a mixture of the degradation effect and the
controller's response.

In Figure~\ref{fig:engine}, note the feedback loop implied by the
closed-loop architecture. The EGT rises as HPC efficiency drops not
just because the thermodynamics change, but because the controller
demands more fuel to compensate. Health monitoring in a closed-loop
system is inherently harder than in an open-loop one.

% -------------------------------------------------------------------
\section{How Degradation Was Represented}
% -------------------------------------------------------------------

\subsection{The two health variables}

The paper represents the health of any engine module using just two
numbers:

\begin{center}
\begin{tabular}{@{}lp{9cm}@{}}
  \toprule
  Variable & Physical meaning \\
  \midrule
  $e(t)$ & Efficiency --- how well the module converts input work into
           useful output. Degradation reduces this. \\
  $f(t)$ & Flow capacity --- how much air mass passes through the module
           per unit time. Degradation generally reduces this too. \\
  \bottomrule
\end{tabular}
\end{center}

These two numbers define the \emph{hidden state} of the engine.
Every sensor output is a function of $(e(t), f(t))$ --- and the
flight conditions.

\subsection{Why this is a reduced-order model}

Real compressor degradation involves thousands of physical variables:
blade surface roughness, tip clearance, seal wear, deposit buildup,
and more. The authors deliberately discard all of this microscopic
detail and retain only two aggregate parameters per module.

This is a \textbf{reduced-order health model}: a deliberate compression
of the full physical state into the minimum number of variables needed
to produce realistic sensor responses. It is not physically perfect.
It does not need to be. Its purpose is to create a
\emph{controllable degradation-to-sensor mapping} --- and for that
purpose, two variables per module are sufficient.

% -------------------------------------------------------------------
\section{Damage Propagation Models}
% -------------------------------------------------------------------

Having chosen C-MAPSS as the simulator, the next question is: how should
$(e(t), f(t))$ evolve over time? The paper reviews three classical
degradation models before arriving at its own.

\subsection{Arrhenius model}

Originally developed for chemical and diffusion failure processes, the
Arrhenius model gives the time to failure as:

\begin{equation}
  t_f = A \, e^{\,\Delta H / kT}
  \label{eq:arrhenius}
\end{equation}

where $T$ is temperature, $k$ is Boltzmann's constant, $A$ is a scaling
factor, and $\Delta H$ is the activation energy. The physical intuition
is simple: higher temperature means more atoms have enough energy to
cross the activation barrier, so degradation processes run faster.

The key insight from Arrhenius: \emph{degradation rate is exponentially
sensitive to temperature}. This is why turbine components --- which run
at extreme temperatures --- degrade faster than compressor components.

\subsection{Coffin-Manson model}

More relevant to mechanical fatigue, the Coffin-Manson model gives the
cycles to failure as:

\begin{equation}
  N_f = A \, f^{-\alpha} \, \Delta T^{-\beta} \, G(T_{\max})
  \label{eq:coffinmanson}
\end{equation}

where $f$ is cycling frequency, $\Delta T$ is the temperature range
per cycle, $G(T_{\max})$ is an Arrhenius term at peak temperature,
and $\alpha$, $\beta$ are exponents. This model captures the damage
accumulated through repeated heating and cooling cycles --- the
dominant failure mode for rotating machinery.

The key insight: \emph{each thermal cycle deposits damage}. The
damage accumulates until the component fails. Number of cycles to
failure depends on both the severity and frequency of the cycling.

\subsection{Eyring model}

The most general of the three, the Eyring model handles multiple
simultaneous stresses:

\begin{equation}
  t_f = A T^{\alpha} \exp\!\left(
    \frac{\Delta H}{kT}
    + \left(B + \frac{C}{T}\right)S_1
    + \left(D + \frac{E}{T}\right)S_2
  \right)
  \label{eq:eyring}
\end{equation}

where $S_1, S_2$ are additional stresses (voltage, pressure, vibration,
etc.) and the constants determine how stresses interact with temperature.
It is the right model when multiple degradation mechanisms act
simultaneously, but it requires many parameters to be calibrated.

\subsection{The common thread: exponential behavior}

All three models share a structural feature: the rate of degradation
\emph{accelerates} over time. Small initial wear creates additional
stress, which causes faster wear, which creates more stress. This
positive feedback loop produces the characteristic exponential-like
growth in damage rate:

\begin{center}
\begin{tabular}{@{}ll@{}}
  small wear & $\to$ more stress \\
  more stress & $\to$ faster wear \\
  faster wear & $\to$ even more stress \\
  \multicolumn{2}{l}{$\to$ runaway degradation} \\
\end{tabular}
\end{center}

\begin{insightbox}
\textbf{The key abstraction:} the authors observe that all three models
produce exponential macro-level behavior, regardless of the microscopic
mechanism. So rather than choosing one physical model and fitting its
parameters, they simply adopt the exponential form directly. The specific
physics is discarded; the correct macroscopic \emph{shape} of degradation
is retained.
\end{insightbox}

\subsection{The chosen model: generalized exponential wear}

Starting from a generalized wear equation $w = Ae^{B(t)}$ with an
upper wear threshold $th_w$, and recasting in terms of health
(normalized remaining safe margin), the paper arrives at:

\begin{equation}
  h(t) = 1 - \exp\!\left\{a t^b\right\}
  \label{eq:health_basic}
\end{equation}

To allow the simulation to start at an arbitrary point in the wear
space (modeling manufacturing variation and prior usage), an initial
deterioration term $d$ is added:

\begin{equation}
  h(t) = 1 - d - \exp\!\left\{a t^b\right\}
  \label{eq:health_full}
\end{equation}

Since both efficiency and flow degrade (with potentially different rates
and shapes), each gets its own version of this equation:

\begin{align}
  e(t) &= 1 - d_e - \exp\!\left\{a_e \, t^{b_e}\right\} \label{eq:eff} \\
  f(t) &= 1 - d_f - \exp\!\left\{a_f \, t^{b_f}\right\} \label{eq:flow}
\end{align}

The parameters of these equations are summarized in
Table~\ref{tab:params}.

\begin{table}[H]
  \centering
  \caption{Parameters of the health equations and their roles.}
  \label{tab:params}
  \begin{tabular}{@{}clp{7.5cm}@{}}
    \toprule
    Parameter & Range & Role \\
    \midrule
    $d$       & $\leq 1\%$ & Initial deterioration. Drawn randomly per engine,
                              models manufacturing variation and prior wear. \\
    $a_k$     & $[0.001, 0.003]$ & Degradation rate. Higher $a$ means faster
                                   deterioration. \\
    $b_k$     & $[1.4, 1.6]$ & Shape exponent. Controls how quickly the
                               exponential accelerates. $b > 1$ produces
                               super-linear growth. \\
    $t$       & $\geq 0$ & Engine cycles. One cycle $\approx$ one flight. \\
    \bottomrule
  \end{tabular}
\end{table}

Figure~\ref{fig:trajectories} shows what this produces: a family of
curves that all start near $h = 1$ and drift toward $h = 0$, but
at different rates and with different shapes depending on their
randomly chosen parameters.

\begin{figure}[H]
  \centering
  \includegraphics[width=\textwidth]{04_degradation_trajectories.png}
  \caption{A family of degradation trajectories $h(t)$ for five engines
  with different parameter sets. All share the same functional form but
  reach failure at different cycle counts. The slow engine lives
  $\sim$180 cycles; the fast one $\sim$100. In the real dataset,
  small process noise makes each trajectory slightly non-monotonic.}
  \label{fig:trajectories}
\end{figure}

% -------------------------------------------------------------------
\section{The Health Index: Defining Failure}
% -------------------------------------------------------------------

\subsection{From module health to overall engine health}

The efficiency and flow trajectories $e(t)$ and $f(t)$ describe the
hidden degradation state. But when does the engine actually fail?
The paper defines failure operationally: an engine fails when it can
no longer operate safely within its design limits under any possible
flight condition.

This is captured through \textbf{operability margins} --- how far the
engine is from violating a safety constraint. Four margins are tracked:

\begin{center}
\begin{tabular}{@{}lp{9cm}@{}}
  \toprule
  Margin & What it measures \\
  \midrule
  $m_\text{Fan}$  & Distance from fan stall. Stall = sudden loss of
                    compression, catastrophic. \\
  $m_\text{LPC}$  & Distance from LPC stall. Same risk. \\
  $m_\text{HPC}$  & Distance from HPC stall. Most relevant given the
                    HPC is the degraded module. \\
  $m_\text{EGT}$  & Distance from maximum exhaust gas temperature.
                    Overtemperature causes turbine blade damage. \\
  \bottomrule
\end{tabular}
\end{center}

Each margin is normalized to $[0, 1]$. The overall health index is the
minimum:

\begin{equation}
  H(t) = \min\!\bigl(m_\text{Fan},\; m_\text{LPC},\;
                     m_\text{HPC},\; m_\text{EGT}\bigr)
  \label{eq:hindex}
\end{equation}

\begin{keybox}
\textbf{Why the minimum?} The engine fails as soon as \emph{any one}
safety boundary is violated. A chain breaks at its weakest link.
$H(t) = 0$ when the first margin is exhausted --- regardless of how
healthy the engine is on all other margins.
\end{keybox}

\subsection{Failure as geometry}

Figure~\ref{fig:geometry} provides the geometric interpretation.
The hidden degradation state is a point in the two-dimensional
$(e\text{ loss},\, f\text{ loss})$ space. As the engine ages,
this point moves away from the healthy origin (bottom-left)
toward the failure region (upper-right).

\begin{figure}[H]
  \centering
  \includegraphics[width=0.90\textwidth]{06_health_index_geometry.png}
  \caption{Degradation as movement through latent health space.
  Contour lines show equal-health-index surfaces. Three trajectories
  with different directions reach the failure boundary at different
  points. The trajectory that moves primarily along the efficiency
  axis reaches failure differently from one that moves primarily
  along the flow axis.}
  \label{fig:geometry}
\end{figure}

The contour lines in Figure~\ref{fig:geometry} represent the health
landscape. Each contour is a locus of $(e, f)$ pairs that produce
the same health index. The failure boundary is the contour at $H = 0$.

The direction of the degradation trajectory matters. An engine that
loses efficiency faster than flow will cross a different boundary
(likely the EGT margin) than one that loses flow faster (likely
the HPC stall margin). This is why the paper discusses Figures~7
and~8 separately in the original: two response surfaces, one
trajectory family.

\subsection{Dependence on operating conditions}

An important subtlety: the margins are not constant. They depend on
flight conditions --- altitude, Mach number, and throttle resolver
angle (TRA). An engine operating at hot-day takeoff faces tighter
margins than the same engine cruising at altitude.

The paper simulated six different flight condition combinations.
The health index computation considers the worst-case condition:
if the engine would violate a margin under \emph{any} simulated
condition, its health index is zero. This ensures the dataset
reflects realistic operational risk.

% -------------------------------------------------------------------
\section{The Application Scenario and Noise Architecture}
% -------------------------------------------------------------------

\subsection{One measurement per flight}

The simulation models each engine flight as a single snapshot taken
at cruise conditions. This reflects real operational practice: health
monitoring systems typically record one representative measurement
per flight, rather than continuous high-rate data.

The effects of between-flight maintenance are \emph{not} explicitly
modeled but are incorporated as process noise --- small random
variations that can occasionally improve performance, making
the efficiency/flow trajectory locally non-monotonic. This is
physically realistic: a cleaning event or minor adjustment can
temporarily recover a few percent of performance.

\subsection{The layered noise architecture}

This is one of the most carefully designed aspects of the dataset.
Four distinct noise sources are stacked, each added at a different
stage of the simulation pipeline. Figure~\ref{fig:noise} shows
the full stack.

\begin{figure}[H]
  \centering
  \includegraphics[width=0.82\textwidth]{05_noise_layer_stack.png}
  \caption{The four-layer noise architecture of the C-MAPSS dataset.
  Each layer introduces a different kind of uncertainty. The layers
  are not independent --- process noise enters the simulation and is
  reshaped by the engine dynamics before reaching the sensor outputs.}
  \label{fig:noise}
\end{figure}

The four layers are:

\textbf{Layer 1 --- Initial wear.} The starting values $e_0$ and $f_0$
are drawn uniformly from $[0.99, 1.0]$. Every simulated engine begins
slightly below perfect health, just as real engines leave the factory
with small imperfections.

\textbf{Layer 2 --- Degradation parameters.} For each engine, the
trajectory parameters $a_k$ and $b_k$ are drawn from a normal
distribution. This produces a \emph{deterministic} trajectory for
each engine --- but each engine has a different one. Given $e_0,
f_0, a_k, b_k$, the health trajectory is fully determined.

\textbf{Layer 3 --- Process noise.} The deterministic trajectory is
then \emph{masked} by a mixture of two normal distributions with
slightly different variances. This is the subtlest noise layer.
A mixture model is harder to characterize than a single Gaussian.
The choice is deliberate: it forces algorithms to deal with
non-trivial uncertainty structure, rather than simple additive
Gaussian noise. This layer also models the between-flight maintenance
effects.

\textbf{Layer 4 --- Measurement noise.} After the noisy health
trajectory is fed through C-MAPSS and steady-state sensor outputs
are computed, a final independent random noise term is added to
each sensor channel. This models sensor imperfection: quantization
error, calibration drift, electrical noise.

\begin{insightbox}
\textbf{Why all four layers?}
Each layer creates a different challenge for the PHM algorithm.
Layer 1 means no two engines start at the same health state.
Layer 2 means no two engines age at the same rate.
Layer 3 means the trajectory cannot be recovered by simple filtering.
Layer 4 means even a perfect model of the engine dynamics would not
give clean health estimates from sensor data.

An algorithm that handles one layer but not the others will fail
in real deployment. The dataset tests all four simultaneously.
\end{insightbox}

% -------------------------------------------------------------------
\section{The Simulation Pipeline}
% -------------------------------------------------------------------

The complete data generation process --- from parameter initialization
to a finished run-to-failure trajectory --- is shown in
Figure~\ref{fig:pipeline}.

\begin{figure}[H]
  \centering
  \includegraphics[width=0.72\textwidth]{08_simulation_pipeline.png}
  \caption{The C-MAPSS simulation loop. Each pass through the main loop
  generates one cycle of sensor data for one engine. The loop continues
  until the health index reaches zero. The entire loop is repeated
  for each engine in the dataset.}
  \label{fig:pipeline}
\end{figure}

The data generation procedure in precise terms:

\begin{enumerate}
  \item Draw $e_0, f_0 \in [0.99, 1.0]$ (initial wear).
  \item Draw $a_k \in [0.001, 0.003]$ and $b_k \in [1.4, 1.6]$ for
        both efficiency and flow (degradation parameters).
  \item Compute the health trajectory $e(t), f(t)$ from
        Equations~\eqref{eq:eff}--\eqref{eq:flow}.
  \item Mask the trajectory with mixture-distribution process noise.
  \item Feed the noisy $(e(t), f(t))$ to C-MAPSS under a randomly
        chosen cruise condition at each step.
  \item Record the 21 sensor outputs once steady state is reached.
  \item Add independent measurement noise to all 21 channels.
  \item Compute the health index $H(t)$ from the four margins.
  \item If $H(t) = 0$: stop. This is the failure point.
  \item Otherwise: increment $t$ and repeat from step~4.
\end{enumerate}

One complete run of this procedure produces one
\emph{run-to-failure trajectory} --- a time series of 21 sensor
readings from some initial cycle count to the failure cycle.

The procedure is repeated hundreds of times to build the full dataset,
with independent parameter draws each time.

% -------------------------------------------------------------------
\section{Competition Dataset Structure}
% -------------------------------------------------------------------

\subsection{What was generated}

The final dataset consists of trajectories simulating HPC degradation
under six different combinations of altitude, TRA, and Mach number.
The six conditions were chosen to span a realistic range of operational
scenarios without being exhaustively comprehensive.

The dataset is divided into three sets with importantly different
properties:

\begin{center}
\begin{tabular}{@{}lp{9.5cm}@{}}
  \toprule
  Set & Properties \\
  \midrule
  Training   & Full trajectories from initial condition to failure.
               The algorithm sees the complete engine life and knows
               when failure occurred. \\
  Test       & Trajectories \emph{truncated} at some point before
               failure. RUL ranges from 10 to 150 cycles.
               Score feedback given after each submission. \\
  Validation & Trajectories truncated before failure. RUL ranges from
               6 to 190 cycles --- a wider distribution than the test
               set. No feedback during competition. \\
  \bottomrule
\end{tabular}
\end{center}

The wider RUL range in the validation set was a hidden robustness test.
Algorithms that overfit to the test set distribution (RUL $\leq 150$)
would fail on validation cases with RUL up to 190. This was not
announced to participants, making the validation set an honest
out-of-distribution test.

\subsection{What participants received and did not receive}

Participants received the 21 sensor readings per cycle for each engine.
They did \emph{not} receive:
\begin{itemize}[noitemsep]
  \item The health index $H(t)$
  \item The four margin values ($m_\text{Fan}$, $m_\text{LPC}$,
        $m_\text{HPC}$, $m_\text{EGT}$)
  \item The degradation parameters $(a_k, b_k, d_e, d_f)$
  \item The hidden state $(e(t), f(t))$
\end{itemize}

All of this had to be inferred from the 21 observable sensor channels.

% -------------------------------------------------------------------
\section{Performance Evaluation and Scoring}
% -------------------------------------------------------------------

\subsection{The asymmetric scoring function}

For any RUL prediction, define the prediction error as:
\[
  d = \hat{t}_\text{RUL} - t_\text{RUL}
  \quad \text{(estimated minus true)}
\]

The aggregate score across $n$ test units is:

\begin{equation}
  s = \begin{cases}
    \displaystyle\sum_{i=1}^{n} e^{-d_i/10} - 1
      & \text{if } d_i < 0 \quad \text{(early prediction)} \\[10pt]
    \displaystyle\sum_{i=1}^{n} e^{\,d_i/13} - 1
      & \text{if } d_i \geq 0 \quad \text{(late prediction)}
  \end{cases}
  \label{eq:score}
\end{equation}

Figure~\ref{fig:scoring} shows this function graphically.

\begin{figure}[H]
  \centering
  \includegraphics[width=0.90\textwidth]{07_asymmetric_scoring_curve.png}
  \caption{The asymmetric scoring function. Both sides grow exponentially
  with error magnitude. The late-prediction side (red, denominator 13)
  is steeper than the early-prediction side (blue, denominator 10),
  encoding the aviation safety preference: it is better to predict
  early and schedule unnecessary maintenance than to predict late and
  risk an in-flight failure. The dotted line shows what a symmetric
  function would look like.}
  \label{fig:scoring}
\end{figure}

\subsection{Why asymmetric?}

The asymmetry encodes a value judgment:

\begin{center}
\begin{tabular}{@{}lll@{}}
  \toprule
  Error type & Consequence & Penalty \\
  \midrule
  Late ($d > 0$) & Engine may fail before maintenance & Severe \\
  Early ($d < 0$) & Unnecessary maintenance cost & Moderate \\
  \bottomrule
\end{tabular}
\end{center}

In aviation, the first consequence is unacceptable. The second is
costly but manageable. The denominators $a_1 = 10$ and $a_2 = 13$
translate this asymmetry into a quantitative penalty structure.

\begin{insightbox}
\textbf{The scoring function is a safety statement encoded as math.}
Choosing $a_1$ and $a_2$ is equivalent to specifying how many
dollars of unnecessary maintenance is worth risking to avoid one
unscheduled engine removal. PHM practitioners should treat these
parameters as domain-specific design choices, not universal constants.
\end{insightbox}

\subsection{The aggregate score and its limitation}

The score in Equation~\eqref{eq:score} aggregates across all test
units using a sum. This creates an important failure mode: two
algorithms with completely different prediction behaviors can produce
identical aggregate scores.

The paper illustrates six prediction scenarios where the aggregate
score is always 6, but the correlation between predicted and true
RUL ranges from 0.0 to 1.0. An algorithm with correlation 0.0 ---
whose predictions have no relationship to ground truth whatsoever ---
can score identically to an algorithm with perfect correlation,
depending on how the errors happen to cancel.

The authors' conclusion: aggregate score should be supplemented by
a correlation metric to ensure \emph{consistent} prediction behavior
across the fleet, not just a lucky average.

% -------------------------------------------------------------------
\section{The Bigger Picture: What This Paper Is}
% -------------------------------------------------------------------

The paper is best understood as three things simultaneously:

\textbf{A methodology paper.} It documents the architecture of a
physics-based degradation simulation pipeline. Every design choice ---
the exponential wear model, the minimum-margin health index, the
layered noise --- is explained and justified. Other researchers can
apply the same architecture to different systems.

\textbf{A dataset paper.} It is the primary citation for the C-MAPSS
dataset, which became one of the most widely used benchmarks in
predictive maintenance research. Understanding the paper is
understanding the dataset's assumptions and limitations.

\textbf{An early digital twin.} The combination of physics-based
simulation + parameterized degradation injection + sensor output
generation is structurally identical to what is now called a
digital twin. The paper was published in 2008. The vocabulary
caught up later.

\begin{keybox}
\textbf{In one sentence:}
The paper answers the question ``how do we generate believable
run-to-failure data for aircraft engines?'' by building a
controllable physics-to-sensor mapping, injecting exponential
degradation into the hidden state, and observing the resulting
sensor trajectories through multiple layers of realistic noise.
\end{keybox}

% -------------------------------------------------------------------
\section*{Appendix: Equation Reference}
% -------------------------------------------------------------------
\addcontentsline{toc}{section}{Appendix: Equation Reference}

\begin{center}
\begin{tabular}{@{}lll@{}}
  \toprule
  Equation & Formula & Where used \\
  \midrule
  Arrhenius & $t_f = Ae^{\Delta H / kT}$ & Background (Sec.~6.1) \\
  Coffin-Manson & $N_f = Af^{-\alpha}\Delta T^{-\beta}G(T_{\max})$ & Background (Sec.~6.2) \\
  Eyring & $t_f = AT^{\alpha}e^{(\cdots)}$ & Background (Sec.~6.3) \\
  Health (basic) & $h(t) = 1 - \exp\{at^b\}$ & Core model \\
  Health (with $d$) & $h(t) = 1 - d - \exp\{at^b\}$ & Actual simulation \\
  Health index & $H(t) = \min(m_\text{Fan}, m_\text{HPC}, m_\text{LPC}, m_\text{EGT})$ & Failure criterion \\
  Score (early) & $s = \sum e^{-d/10} - 1$ & Competition metric \\
  Score (late) & $s = \sum e^{d/13} - 1$ & Competition metric \\
  \bottomrule
\end{tabular}
\end{center}

% -------------------------------------------------------------------
\end{document}
